\section{Requisitos e geração do arquivo final (PDF)}

A versão final da dissertação ou tese será entregue \textbf{exclusivamente em formato digital (PDF)}, composta por um \textbf{PDF mestre} contendo todo o conteúdo do trabalho e os documentos obrigatórios. Quando aplicável, também deve ser entregue uma \textbf{cópia protegida} (Seção~\ref{sec:cap6-protecao}).

\subsection{Preparação do documento-fonte (recomendado)}

\begin{enumerate}
    \item Antes da geração do PDF, o arquivo de origem (Word/LibreOffice/\LaTeX{} etc.) deve estar \textbf{estruturado com níveis de títulos} (ex.: Título 1, Título 2, Título 3\ldots), de modo a permitir a \textbf{criação automática de marcadores (bookmarks)}.
    \item Atualize o \textbf{Sumário} no documento-fonte e revise se todas as seções/subseções estão corretamente listadas.
\end{enumerate}

\subsection{Geração do PDF mestre}

\begin{enumerate}
    \item Gere o PDF diretamente a partir do documento-fonte, preferencialmente usando uma opção de exportação que \textbf{crie marcadores a partir dos títulos/estrutura}.

    \item Converta e incorpore ao PDF mestre:
    \begin{itemize}
        \item a \textbf{Folha de Registro do Documento} (quando fornecida em arquivo digital);
        \item eventuais páginas/elementos obrigatórios (pré-textuais e pós-textuais).
    \end{itemize}

    \item \textbf{Anexos não disponíveis em meio digital} devem ser digitalizados e inseridos no PDF mestre. Para anexos textuais, recomenda-se aplicar \ac{OCR} para permitir pesquisa e seleção de texto, sempre que tecnicamente viável.

    \item Substitua a \textbf{Folha de Rosto} (ou página equivalente, quando aplicável) pela \textbf{versão assinada} (escaneada com qualidade adequada ou assinada digitalmente conforme regras institucionais).

    \item Ao final, revise a integridade do PDF mestre:
    \begin{itemize}
        \item ordem das páginas;
        \item qualidade de leitura de páginas escaneadas (sem cortes, orientação correta);
        \item consistência do Sumário e dos marcadores (ver Seção~\ref{sec:cap6-bookmarks});
        \item metadados básicos (título, autor, ano), quando aplicável.
    \end{itemize}
\end{enumerate}

\subsection{Ferramentas recomendadas}

São recomendadas ferramentas que permitam gerar e editar PDF \textbf{localmente} (sem \textit{upload}), sobretudo quando o documento contiver dados pessoais ou informações sensíveis. A Tabela~\ref{tab:ferramentas-pdf} lista opções de \ac{FOSS} e freeware, com interfaces \ac{GUI} ou \ac{CLI}.

\begin{table}[H]
  \centering
  \caption{Ferramentas recomendadas para manipulação de PDF.}
  \label{tab:ferramentas-pdf}
  \begin{tabular}{lllll}
    \toprule
    \textbf{Ferramenta} & \textbf{Licença} & \textbf{Interface} & \textbf{Plataforma} & \textbf{Uso} \\
    \midrule
    PDF24 Creator & Freeware & \ac{GUI} & Windows & Edição, mesclagem \\
    Stirling PDF & \ac{FOSS} & \ac{GUI} & Multi & Edição geral \\
    \texttt{qpdf} & \ac{FOSS} & \ac{CLI} & Multi & Criptografia \\
    \texttt{pdftk} & \ac{FOSS} & \ac{CLI} & Multi & Mesclagem \\
    JSignPdf & \ac{FOSS} & \ac{GUI} & Multi & Assinatura digital \\
    \bottomrule
  \end{tabular}
\end{table}

\subsection{Nome do arquivo}

O arquivo final deve ser nomeado como:

\begin{center}
\texttt{DM/TD999\_YYYY.pdf}
\end{center}

onde:

\begin{itemize}
    \item \textbf{DM} = Dissertação de Mestrado; \textbf{TD} = Tese de Doutorado;
    \item \textbf{999} = número de registro atribuído;
    \item \textbf{YYYY} = ano da defesa.
\end{itemize}

\section{Marcadores (bookmarks) no PDF}\label{sec:cap6-bookmarks}

\subsection{Requisitos mínimos de marcadores}

No PDF mestre devem existir marcadores para:

\begin{enumerate}
    \item cada elemento \textbf{pré-textual} (ex.: folha de rosto, resumo, \textit{abstract}, listas etc., conforme aplicável);
    \item o \textbf{Sumário}, incluindo \textbf{todas as seções e subseções} nele indicadas;
    \item a \textbf{Folha de Registro do Documento};
    \item \textbf{Anexos/Apêndices}, quando existirem.
\end{enumerate}

\subsection{Método preferencial (automático e reprodutível)}

\begin{enumerate}
    \item Estruture o documento-fonte com estilos de título (níveis hierárquicos).
    \item Em \LaTeX{}, use o pacote \texttt{hyperref} para gerar marcadores automaticamente (ex.: inclua \texttt{hyperref} no preâmbulo) e garanta que as seções/subseções estejam bem estruturadas.
    \item Ao exportar para PDF, habilite a opção equivalente a ``\textbf{criar marcadores a partir de títulos/estrutura}''.
    \item Após exportar, \textbf{valide}:
    \begin{itemize}
        \item se todos os itens do Sumário possuem marcador correspondente;
        \item se a hierarquia (pai/filho) reflete o nível de seção/subseção;
        \item se cada marcador aponta para a página correta.
    \end{itemize}
\end{enumerate}

\subsection{Ajuste manual (quando necessário)}

Quando a exportação automática não produzir marcadores adequados, use um editor de PDF para:

\begin{itemize}
    \item renomear marcadores;
    \item corrigir destinos;
    \item ajustar a hierarquia (aninhamento);
    \item inserir marcadores faltantes.
\end{itemize}

\noindent\textbf{Critério de aceitação:} ao clicar em cada marcador, o leitor deve ser direcionado imediatamente ao trecho correspondente, com navegação clara.

\section{Proteção do documento e cópias de distribuição}\label{sec:cap6-protecao}

\subsection{Princípios}

\begin{enumerate}
    \item A proteção pode ter finalidades distintas:
    \begin{itemize}
        \item \textbf{Confidencialidade}: exigir senha para abrir o PDF;
        \item \textbf{Restrição de permissões}: desencorajar edição/impressão/cópia (sem impedir absolutamente);
        \item \textbf{Autenticidade/integridade}: assinatura digital para evidenciar alterações.
    \end{itemize}

    \item Quando houver exigência de compatibilidade com \textbf{PDF/A} ou com repositórios que rejeitam criptografia, recomenda-se:
    \begin{itemize}
        \item manter o \textbf{PDF mestre sem criptografia}; e
        \item gerar \textbf{uma cópia protegida separada}, se a norma exigir restrição.
    \end{itemize}
\end{enumerate}

\subsection{Ferramentas recomendadas (\ac{FOSS}/freeware)}

\begin{itemize}
    \item \textbf{\texttt{qpdf} (\ac{FOSS}, \ac{CLI})}: criptografia AES-256 e permissões; ideal para \textit{workflow} reprodutível.
    \item \textbf{\texttt{pdftk} / \texttt{pdftk-java} (\ac{FOSS}/free, \ac{CLI})}: permissões/criptografia simples.
    \item \textbf{PDF24 Creator (freeware, \ac{GUI} Windows)}: aplicar senha/permissões e montar/mesclar PDFs localmente.
    \item \textbf{Stirling PDF (\ac{FOSS})}: operações de PDF via interface local.
    \item \textbf{JSignPdf (freeware/\ac{FOSS})}: assinatura digital.
\end{itemize}

\subsection{Exemplos por linha de comando}

\subsubsection*{Criptografar exigindo senha para abrir (confidencialidade)}

\textbf{qpdf (AES-256):}

\begin{Verbatim}[breaklines, breakanywhere, numbers=left]
qpdf --encrypt "SENHA_ABRIR" "SENHA_DONO" 256 -- input.pdf output_protegido.pdf
\end{Verbatim}

\subsubsection*{Restringir impressão/edição/cópia (permissões)}

\textbf{qpdf (sem senha para abrir; com senha de permissões):}

\begin{Verbatim}[breaklines, breakanywhere, numbers=left]
qpdf --encrypt "" "SENHA_DONO" 256 --print=none --modify=none --extract=n --annotate=n --assemble=n --form=n -- input.pdf output_restrito.pdf
\end{Verbatim}

\textbf{pdftk (exemplo de permissões):}

\begin{Verbatim}[breaklines, breakanywhere, numbers=left]
pdftk input.pdf output restrito.pdf owner_pw SENHA_DONO user_pw "" allow DegradedPrinting
\end{Verbatim}

\subsubsection*{\ac{OCR} em anexos escaneados (para PDF pesquisável)}

Se você estiver no Windows e preferir usar a suíte do PDF24:

\begin{Verbatim}[breaklines, breakanywhere, numbers=left]
pdf24-Ocr.exe -outputFile saida.pdf -language por -dpi 300 "entrada.pdf"
\end{Verbatim}

\noindent (Alternativamente, é possível aplicar \ac{OCR} via outras ferramentas locais conforme disponibilidade institucional.)

\subsection{Checklist de entrega (sugerido)}

\begin{itemize}
    \item[$\Box$] PDF mestre completo (todas as páginas e anexos) e na ordem final
    \item[$\Box$] Folha assinada corretamente inserida
    \item[$\Box$] Marcadores completos e testados (pré-textuais, Sumário com seções e subseções, Folha de Registro)
    \item[$\Box$] Anexos escaneados legíveis; \ac{OCR} aplicado quando praticável
    \item[$\Box$] Arquivo nomeado como \texttt{DM/TD999\_YYYY.pdf}
    \item[$\Box$] (Se exigido) cópia protegida gerada e verificada em um leitor de PDF
\end{itemize}